% Generated by reproduce/make_research_tables.py; edit the generator, not this file.
\begin{table}[H]
\centering\footnotesize
\setlength{\tabcolsep}{3pt}
\renewcommand{\arraystretch}{1.12}
\begin{tabular*}{\linewidth}{@{\extracolsep{\fill}}lrrrrr@{}}
\toprule
Case & $n$ & \shortstack{Faithful\\scalar bits} & \shortstack{Laurent\\coefficient bits} & \shortstack{Global\\integer bits} & \shortstack{Normalized\\mantissa bits}\\
\midrule
Tree, depth 6 & 126 & 63,757 & 2 & 32,321 & 768\\
Tree, depth 7 & 254 & 258,573 & 2 & 130,177 & 1,536\\
Grid $10\times10$ & 100 & 10,855 & 2 & 10,921 & 832\\
\bottomrule
\end{tabular*}
\caption{Maximum bit length of one stored scalar, Laurent coefficient, or normalized mantissa in completed queries from the final paired run. The values agree across all repeated samples. They are not total table-memory measurements: the normalized mode also stores a valuation, and every mode stores frontier keys and other metadata. The two integer tensor modes use the same exact polynomial model and certified order.}
\label{tab:tensor-precision}
\end{table}
